% ==============================================================================
% front/introduction.tex -- Introduction (original book pages 73-76)
% ==============================================================================
\thispagestyle{plain}

This expose represents an attempt to understand some of the recent work of
Atkin, Swinnerton-Dyer, and Serre on the \uline{congruence properties} of the
q-expansion coefficients of modular forms\index{modular form} from the point of view of the theory
of moduli of elliptic curves\index{elliptic curve}, as developed abstractly by Igusa and recently
reconsidered by Deligne. In this optic, a modular form of weight $k$ and level
$n$ becomes a section of a certain line bundle $\w^{\tens k}$ on the modular
variety $M_n$ which "classifies" elliptic curves with level $n$ structure (the
level $n$ structure is introduced for purely technical reasons). The modular
variety $M_n$ is a smooth curve over $\Z[1/n]$, whose "physical appearance" is
the same whether we view it over $\C$ (where it becomes $\varphi(n)$ copies of
the quotient of the upper half plane by the principal congruence subgroup
$\Gamma(n)$ of $SL(2,\Z)$) or over the algebraic closure of $\Z/p\Z$, (by
"reduction modulo $p$") for primes $p$ not dividing $n$.

\uline{This very fact rules out} the possibility of obtaining p-adic properties
of modular forms \uline{simply} by studying the geometry of
$M_n\tens\Z/p\Z$ and its line bundles $\w^{\tens k}$; we can only obtain the
reductions modulo $p$ of identical relations which hold over $\C$.

The key is instead to \uline{isolate} the finite set of points of
$M_n\tens\Z/p\Z$ corresponding to \uline{supersingular} elliptic curves in
characteristic $p$, those whose \uline{Hasse invariant\index{Hasse invariant} vanishes}. One then
considers various "rigid-analytic" open subsets of $M_n\tens\Z_p$ defined by
\uline{removing} p-adic discs of various radii around the supersingular points
in characteristic $p$. This makes sense because the Hasse invariant is the
reduction modulo $p$ of a true modular form (namely $E_{p-1}$) over $\Z_p$, so
we can define a rigid analytic\index{rigid analytic space} open subset of $M_n\tens\Z_p$ by taking only
those p-adic elliptic curves on which $E_{p-1}$ has p-adic absolute value
greater than some $\varepsilon>0$. We may then define various sorts of truly
p-adic modular forms\index{p-adic modular form} as functions of elliptic curves on which
$|E_{p-1}|>\varepsilon$, or equivalently as sections of the line bundles
$\w^{\tens k}$ \uline{restricted} to the above-constructed rigid analytic open
sets of $M_n\tens\Z_p$. [The role of the choice of $\varepsilon$ is to specify
the rate of growth of the coefficients of the Laurent series development
around the "missing" supersingular points].

The most important tool in the study of these p-adic modular forms is the
endomorphism they undergo by a "canonical lifting of the Frobenius\index{Frobenius}
endomorphism" from characteristic $p$. This endomorphism comes about as
follows. Any elliptic curve on which $|E_{p-1}|>\varepsilon$ for suitable
$\varepsilon$ carries a "\uline{canonical subgroup\index{canonical subgroup}}" of order $p$, whose
reduction modulo $p$ is the Kernel of Frobenius. The "canonical lifting" above
is the \uline{endomorphism} obtained by \uline{dividing} the universal elliptic
curve by \uline{its} canonical subgroup (over the rigid open set of
$M_n\tens\Z_p$ where it exists).

This endomorphism is related closely to Atkin's work. His operator $U$ is
simply ($\tfrac1p$ times) the \uline{trace} of the canonical lifting of
Frobenius, and certain of his results on the q-expansion of the function $j$
may be interpreted as statements about the \uline{spectral theory} of the
operator $U$.

The relation to the work of Swinnerton-Dyer and Serre is more subtle, and
depends on the fact that the data of the action of the "canonical lifting of
Frobenius" on $\w^{-1}$ over the rigid open set $|E_{p-1}|\geq 1$ is
\uline{equivalent} to the knowledge of the representation of the
\uline{fundamental group} of the open set of $M_n\tens\Z/p\Z$ where the Hasse
invariant is invertible on the \uline{p-adic Tate module} $T_p$ (which for a
non-supersingular curve in characteristic $p$ is a free $\Z_p$-module of rank
one). Thanks to Igusa, we know that this representation is as
\uline{non-trivial} as possible, and this fact, \uline{interpreted in terms of
the action} of the canonical Frobenius on the $\w^{\tens k}$, leads to certain
of the congruences of Swinnerton-Dyer and Serre.

In the first chapter, we review without proof certain aspects of the moduli of
elliptic curves, and deduce various forms of the "q-expansion principle\index{q-expansion principle}."
This chapter owes much (probably its very existence) to discussions with
Deligne. It is not "p-adic", and may be read more or less independently of the
rest of the paper.

The second chapter develops at length various "p-adic" notions of modular
form, in the spirit described above. A large part of it ($r\neq 1$) was
included with an eye to Dwork-style applications to Atkin's work, and may be
omitted by the reader interested only in Swinnerton-Dyer and Serre style
congruences. The \uline{idea} of working at such "p-adic modular forms" is due
entirely to Serre, who in his 1972 College de France course stressed their
importance.

The third chapter develops the theory of the "canonical subgroup." This theory
is due \uline{entirely} to Lubin, who has unfortunately not published it except
for a tiny hint \cite{lubin-newton}. The second half of the chapter interprets certain
congruences of Atkin in terms of p-adic Banach\index{Banach space} spaces, the spectrum of the
operator $U$, etc. The possibility of this interpretation is due to Dwork,
through his realization that not only is $pU$ integral, but $U$ itself is
"essentially" integral (cf\cite{dwork-involume}).

The fourth chapter explains the relation between the canonical Frobenius and
certain congruences of Swinnerton-Dyer and Serre. It begins by recalling a
"coherent sheaf" description of p-adic representations of the fundamental
group of certain schemes on which $p$ is nilpotent. This description is
certainly well-known, and basically due to Hasse and Witt, but does not seem to
be recorded elsewhere in the form we require. Using it, we show that the
representation corresponding to $\w$ with its canonical Frobenius is that
afforded by the (rank-one) p-adic Tate module $T_p$ of non-supersingular
elliptic curves. We then prove the extreme \uline{non-triviality} of this
representation in "canonical subgroup" style. This non-triviality is due to
Igusa, whose proof is finally not so different from the one given. We then
apply this result of non-triviality to deduce certain of the congruences of
Swinnerton-Dyer and Serre.

In the first appendix, which is a sort of "chapter zero", we explain the
relation between the classical approach to elliptic curves via their period
lattices and the "modern" one, the relation of DeRham cohomology of elliptic
curves to modular forms, and the relation between the Gauss-Manin connection\index{Gauss--Manin connection},
Ramanujan's function $P(q)$, and Serre's $\partial$-operator on modular forms.
The results are due to Weierstrass and Deligne. It is concluded by a "table" of
formulas.

The second appendix explains the relation between the canonical Frobenius on
p-adic modular forms and the Frobenius endomorphism of the DeRham cohomology
of elliptic curves. It may also be read as an appendix to \cite{katz73}.

The third appendix relates Hecke polynomials mod $p$ to L-series, coherent
cohomology and the Fredholm determinant\index{Fredholm determinant} of $U$.

As should by now be obvious, this expose owes its very existence to Lubin,
Serre, Deligne, Atkin, and Dwork. It is a pleasure to acknowledge my debt to
them, and to thank M. Rapoport for many helpful discussions.
